\documentclass[10pt,aspectratio=169]{beamer}
% Preámbulo modular para presentaciones Beamer (Modelación Estadística)
% Diseño y paleta institucional del Tecnológico de Monterrey

\documentclass[aspectratio=169,xcolor={dvipsnames,table}]{beamer}

\usepackage[utf8]{inputenc}
\usepackage[spanish,mexico]{babel}
\usepackage{amsmath,amssymb,amsfonts,mathtools}
\usepackage{graphicx}
\usepackage{listings}
\usepackage{booktabs}
\usepackage{tikz}
\usetikzlibrary{matrix,arrows,decorations.pathmorphing,shapes.geometric,positioning}

% Paleta Institucional Tecnológico de Monterrey
\definecolor{TecAzul}{HTML}{0039A6}       % Azul TEC: color primario institucional
\definecolor{TecAzulOscuro}{HTML}{1A2E51} % Azul marino oscuro
\definecolor{TecRojo}{HTML}{EC2661}       % Rojo de acento (paleta miTec)
\definecolor{TecGris}{HTML}{646464}       % Gris neutro para comentarios y subtítulos
\definecolor{TecFondoBloque}{HTML}{F4F6F9}% Fondo suave para bloques

% Configuración de Tema Beamer
\usetheme{Boadilla}
\usecolortheme[named=TecAzulOscuro]{structure}
\setbeamercolor{alerted text}{fg=TecRojo}
\setbeamercolor{palette primary}{bg=TecAzulOscuro,fg=white}
\setbeamercolor{palette secondary}{bg=TecAzul,fg=white}
\setbeamercolor{palette tertiary}{bg=TecAzulOscuro!80!black,fg=white}
\setbeamercolor{title}{fg=TecAzulOscuro,bg=TecFondoBloque}
\setbeamercolor{frametitle}{fg=TecAzulOscuro,bg=TecFondoBloque}

% Bloques personalizados elegantes
\setbeamercolor{block title}{bg=TecAzulOscuro,fg=white}
\setbeamercolor{block body}{bg=TecFondoBloque,fg=black}
\setbeamercolor{block title alerted}{bg=TecRojo,fg=white}
\setbeamercolor{block body alerted}{bg=TecRojo!8,fg=black}
\setbeamercolor{block title example}{bg=TecAzul,fg=white}
\setbeamercolor{block body example}{bg=TecAzul!8,fg=black}

% Redondear bordes y añadir sombra sutil
\setbeamertemplate{blocks}[rounded][shadow=true]
\setbeamertemplate{navigation symbols}{}

% Pie de página limpio con número de diapositiva
\setbeamertemplate{footline}{
  \leavevmode%
  \hbox{%
  \begin{beamercolorbox}[wd=.7\paperwidth,ht=2.25ex,dp=1ex,left]{author in head/foot}%
    \hspace*{2ex}\insertshortauthor~~(\insertshortinstitute) \hfill \insertshorttitle
  \end{beamercolorbox}%
  \begin{beamercolorbox}[wd=.3\paperwidth,ht=2.25ex,dp=1ex,right]{date in head/foot}%
    \insertshortdate{}\hspace*{2em}
    \insertframenumber{} / \inserttotalframenumber\hspace*{2ex}
  \end{beamercolorbox}}%
  \vskip0pt%
}

% Configuración de Listings de Python para visualización pedagógica de código
\lstset{
    language=Python,
    basicstyle=\ttfamily\footnotesize,
    keywordstyle=\color{TecAzulOscuro}\bfseries,
    stringstyle=\color{TecRojo},
    commentstyle=\color{TecGris}\itshape,
    showstringspaces=false,
    breaklines=true,
    frame=lines,
    rulecolor=\color{TecAzulOscuro!30},
    backgroundcolor=\color{TecFondoBloque},
    aboveskip=0.5em,
    belowskip=0.5em,
    literate={á}{{\'a}}1 {é}{{\'e}}1 {í}{{\'i}}1 {ó}{{\'o}}1 {ú}{{\'u}}1
             {Á}{{\'A}}1 {É}{{\'E}}1 {Í}{{\'I}}1 {Ó}{{\'O}}1 {Ú}{{\'U}}1
             {ñ}{{\~n}}1 {Ñ}{{\~N}}1 {¿}{{?`}}1 {¡}{{!`}}1
}

% Metadatos institucionales por defecto
\title[Modelación Estadística]{Modelación Estadística para Ciencia de Datos e Ingeniería}
\author[J. Castillo Colmenares]{Juliho David Castillo Colmenares}
\institute[Tec de Monterrey]{Tecnológico de Monterrey \\ Escuela de Ingeniería y Ciencias}
\date{\today}
% Comandos y macros estadísticas compartidas para las presentaciones Beamer (Metropolis)

% Conjuntos numéricos básicos
\newcommand{\R}{\mathbb{R}}
\newcommand{\N}{\mathbb{N}}
\newcommand{\Q}{\mathbb{Q}}
\newcommand{\Z}{\mathbb{Z}}
\newcommand{\set}[1]{\left\{ #1 \right\}}
\newcommand{\abs}[1]{\left|#1\right|}
\newcommand{\norm}[1]{\left\| #1 \right\|}

% Comandos estadísticos y probabilísticos del curso
\newcommand{\Var}{\operatorname{Var}}
\newcommand{\cov}{\operatorname{Cov}}
\newcommand{\corr}{\rho}
\newcommand{\comb}[2]{\begin{pmatrix} #1 \\ #2 \end{pmatrix}}
\newcommand{\s}{\sigma}
\newcommand{\particion}{\mathfrak{P}}
\newcommand{\card}[1]{\#\left( #1 \right)}
\newcommand{\seq}[3]{\set{#1}_{#2}^{#3}}
\newcommand{\E}{\mathbb{E}}
\newcommand{\Prob}{\mathbb{P}}

% Entornos pedagógicos y cajas en español académico natural
\newenvironment{discusion}{%
  \begin{alertblock}{Pregunta de Discusión}%
}{%
  \end{alertblock}%
}

\newenvironment{reflexion}{%
  \begin{alertblock}{Reflexión para el Aula}%
}{%
  \end{alertblock}%
}

\newenvironment{paradoja}{%
  \begin{alertblock}{Paradoja de Exploración}%
}{%
  \end{alertblock}%
}

\newenvironment{motivacion}{%
  \begin{exampleblock}{Motivación: Ciencia de Datos en la Práctica}%
}{%
  \end{exampleblock}%
}

\newenvironment{retoclase}{%
  \begin{block}{Reto Rápido en Equipo}%
}{%
  \end{block}%
}

\title[Comparaciones post-hoc]{Sección 08.04: Comparaciones múltiples y pruebas post-hoc}\subtitle{Después de un ANOVA significativo}
\author[J. Castillo Colmenares]{Juliho Castillo Colmenares}\institute{Tecnológico de Monterrey}\date{}
\begin{document}
\begin{frame}[plain]\titlepage\end{frame}
\begin{frame}{Hoja de ruta}\small\begin{enumerate}\item Por qué comparar después de ANOVA.\item LSD, Tukey y Tukey--Kramer.\item Bonferroni y reporte.\end{enumerate}\end{frame}
\begin{frame}{Fuente y alcance}\small La sección parte de un rechazo global de ANOVA y localiza las diferencias entre tratamientos sin ignorar el error familiar.\begin{block}{Pregunta guía}¿Qué pares de medias explican el rechazo global?\end{block}\end{frame}
\begin{frame}{Límite del contraste global}\small ANOVA solo concluye que al menos una media difiere. Comparar todos los pares sin ajuste eleva la probabilidad de falsos positivos.\end{frame}
\begin{frame}{Diferencia mínima de Fisher}\small Para dos medias, LSD usa el cuadrado medio del error del ANOVA:\[\mathrm{LSD}=t_{\alpha/2,N-k}\sqrt{\mathrm{CME}\left(\frac1{n_i}+\frac1{n_j}\right)}.\]\end{frame}
\begin{frame}{Criterio LSD}\small Declarar diferencia cuando\[|\bar Y_{i.}-\bar Y_{j.}|>\mathrm{LSD}.\]Es potente, pero no controla por sí sola el error familiar con muchos tratamientos.\end{frame}
\begin{frame}{LSD protegido}\small Una práctica prudente calcula LSD solo después de un $F$ global significativo. Aun así, con muchos grupos puede quedar inflación residual.\end{frame}
\begin{frame}{Tukey HSD}\small Tukey compara todos los pares y controla el error familiar. Para tamaños iguales,\[\mathrm{HSD}=q_{\alpha,k,N-k}\sqrt{\frac{\mathrm{CME}}n}.\]\end{frame}
\begin{frame}{Criterio Tukey}\small Para cada par, declarar diferencia si\[|\bar Y_{i.}-\bar Y_{j.}|>\mathrm{HSD}.\]El rango estudentizado incorpora el número de tratamientos.\end{frame}
\begin{frame}{Tukey--Kramer}\small Con tamaños desiguales,\[\mathrm{HSD}_{ij}=\frac{q_{\alpha,k,N-k}}{\sqrt2}\sqrt{\mathrm{CME}\left(\frac1{n_i}+\frac1{n_j}\right)}.\]\end{frame}
\begin{frame}{Corrección de Bonferroni}\small Si hay $m=\binom{k}{2}$ comparaciones, usar\[\alpha^*=\frac{\alpha}{m}.\]Es universal y conservadora.\end{frame}
\begin{frame}{Procedimiento I}\small\begin{enumerate}\item Confirmar el $F$ global.\item Enumerar los pares.\item Elegir LSD, Tukey o Bonferroni según el objetivo.\end{enumerate}\end{frame}
\begin{frame}{Procedimiento II}\small\begin{enumerate}\setcounter{enumi}{3}\item Calcular el umbral de cada par.\item Comparar diferencias absolutas.\item Ajustar el lenguaje al método.\end{enumerate}\end{frame}
\begin{frame}{Aplicación: medias}\small Suponga $k$ tratamientos con medias $\bar Y_{i.}$, tamaños $n_i$ y CME del ANOVA. Las comparaciones usan la misma estimación global del error.\end{frame}
\begin{frame}{Aplicación: LSD}\small Calcule cada LSD con $t_{\alpha/2,N-k}$. Una diferencia que excede el umbral se declara detectable bajo el procedimiento elegido.\end{frame}
\begin{frame}{Aplicación: Tukey}\small Use el cuantil $q$ para controlar simultáneamente todos los pares. La conclusión es más conservadora que LSD.\end{frame}
\begin{frame}{Aplicación: tamaños desiguales}\small Si $n_i\neq n_j$, sustituya HSD por Tukey--Kramer y conserve los grados de libertad del error.\end{frame}
\begin{frame}{Aplicación: Bonferroni}\small Divida $\alpha$ entre el número de pares. El umbral aumenta, reduciendo falsos positivos y también potencia.\end{frame}
\begin{frame}{Interpretación}\small Una comparación ajustada significativa identifica un par compatible con una diferencia; no convierte todas las diferencias numéricas en evidencia estadística.\end{frame}
\begin{frame}{Reporte}\small Incluya el $F$ global, método post-hoc, umbrales, pares detectados, estimaciones de diferencia e intervalos cuando estén disponibles.\end{frame}
\begin{frame}{Síntesis}\small\begin{itemize}\item ANOVA responde una pregunta global.\item LSD prioriza potencia.\item Tukey controla pares simultáneos.\item Bonferroni es simple y conservadora.\end{itemize}\end{frame}
\begin{frame}{Conexión con el capítulo}\small Las decisiones post-hoc deben respetar el diseño, sus supuestos y la pregunta científica; no son una etapa automática separada del experimento.\end{frame}
\end{document}
